% -----------------------------------------------
% Template for ISMIR LBD Papers
% 2025 version, based on previous ISMIR templates

% Requirements :
% * 2+1 page length maximum
% * 10MB maximum file size
% * Copyright note must appear in the bottom left corner of first page
% * Clearer statement about citing own work in anonymized submission
% (see conference website for additional details)
% -----------------------------------------------

\documentclass{article}
\usepackage[T1]{fontenc} % add special characters (e.g., umlaute)
\usepackage[utf8]{inputenc} % set utf-8 as default input encoding
\usepackage{amsmath,cite,url}
\usepackage{graphicx}
\usepackage[lbd]{ismir}

\usepackage{color}

\def\lbd{}	% Flag to use correct LBD settings in the paper, please do not modify this line


% \onecolumn
% Title. Please use IEEE-compliant title case when specifying the title here,
% as it has implications for the copyright notice
% ------
\title{LayerIt: A Python Framework for Score-Aligned MIR Visualizations}

% Note: Please do NOT use \thanks or a \footnote in any of the author markup

% Single address
% To use with only one author or several with the same address
% ---------------
%\oneauthor
% {Names}
% {Affiliations}

% Two addresses
% --------------
%\twoauthors
%  {First author} {School \\ Department}
%  {Second author} {Company \\ Address}

% Three addresses
% --------------
\threeauthors
  {First Author} {Affiliation1 \\ {\tt author1@ismir.edu}}
  {Second Author} {Affiliation2 \\ {\tt author2@ismir.edu}}
  {Third Author} {Affiliation3 \\ {\tt author3@ismir.edu}}

% Four or more addresses
% OR alternative format for large number of co-authors
% ------------
%\multauthor
%{First author$^1$ \hspace{1cm} Second author$^1$ \hspace{1cm} Third author$^2$} { \bfseries{Fourth author$^3$ \hspace{1cm} Fifth author$^2$ \hspace{1cm} Sixth author$^1$}\\
%  $^1$ Department of Computer Science, University , Country\\
%$^2$ International Laboratories, City, Country\\
%$^3$  Company, Address\\
%{\tt\small CorrespondenceAuthor@ismir.edu, PossibleOtherAuthor@ismir.edu}
%}

% For the author list in the Creative Common license and PDF metadata, please enter author names.
% Please abbreviate the first names of authors and add 'and' between the second to last and last authors.
\def\authorname{F. Author, S. Author, and T. Author}

% Optional: To use hyperref, uncomment the following.
%\usepackage[bookmarks=false,pdfauthor={\authorname},pdfsubject={\papersubject},hidelinks]{hyperref}
% Mind the bookmarks=false option; bookmarks are incompatible with ismir.sty.

\sloppy % please retain sloppy command for improved formatting

\begin{document}

%
\maketitle
%
\begin{abstract}
MIR visualization tools play a crucial role in music analysis, annotation, and evaluation, but existing solutions often make it difficult to combine audio, symbolic scores, and algorithmic outputs in a single, time-aligned view. This thesis presents LayerIt!, a Python-based object-oriented framework for generating custom score-aligned MIR visualizations by combining established libraries and tools, including librosa, matplotlib, Modusa, ScoreWarp, and BeatThis. LayerIt is built around a modular layer-composition model in which data sources, audio plots, score renderings, and algorithmic outputs are treated as independent components that can be assembled into tailored visualizations. The resulting figures are exported as SVG, enabling resolution-independent output, post-processing, and future integration with web-based or interactive applications. As a validation case, LayerIt is used to visualize beat tracking results alongside aligned score and audio representations, showing how the framework supports both beat annotation and beat tracking evaluation through discrete beat estimates and continuous confidence curves. More broadly, the project demonstrates a reusable foundation for future specialized MIR visualization workflows.
\end{abstract}
%
\section{Introduction}\label{sec:introduction}
MIR increasingly depends on visualizations that combine multiple modalities in a single time reference. Researchers often need to inspect symbolic scores, audio representations, and algorithmic output together, but existing workflows typically rely on \textit{ad hoc} image editing or GUI tools that are not designed for flexible figure composition. This is especially visible in audio-to-score alignment, where the value of the alignment is not only the algorithmic result itself, but also the ability to present it in a way that is readable, reproducible, and easy to revise.

This Late-Breaking Demo presents LayerIt!, the visualization framework developed in my thesis. The project grew out of a broader investigation into beat tracking and beat annotation, but the final contribution is not a new alignment algorithm. Instead, it is a reusable visualization layer that makes score-aligned MIR figures easier to build and extend. The framework is meant to complement existing tools such as Piano Precision, which demonstrated the usefulness of score-aware GUI workflows, while shifting the emphasis toward programmable, publication-ready visual composition \cite{jiangPIANOPRECISIONADVANCING2025}.

\section{Motivation}
MIR research increasingly depends on visualizations that integrate multiple modalities in a single time reference. Researchers often need to inspect algorithmic outputs alongside audio representations and symbolic scores in a time-coherent manner, and this kind of multi-modal visualization is central to both interpretation and analysis. Audio-to-Score Alignment provides the temporal anchor that makes this possible by synchronizing a performance recording with its corresponding symbolic score, allowing heterogeneous data sources to be aligned into one unified view.

The challenge is that these visualizations are often assembled through manual editing or tool-specific workarounds. Static figures can be informative, but they are difficult to revise when the underlying data changes, and they rarely preserve a clear relationship between the different visual layers. Existing GUI tools are useful for inspection, yet they are not designed for flexible composition of score-aligned views that combine plots, score renderings, and derived data in a reusable way. As a result, researchers frequently rely on 	extit{ad hoc} methods that are hard to reproduce and hard to adapt.

LayerIt was developed to address this gap by making score-aligned MIR visualizations more composable, reusable, and easier to adapt to changing data and research needs. The goal is not to replace the underlying alignment or analysis methods, but to provide a visual framework that lets those methods be presented in a clearer and more maintainable way.

\section{LayerIt Framework}

LayerIt is an object-oriented Python framework that treats each visualization component as a layer with a common interface. A layer can load its own data, draw itself, and be combined with other layers into a higher-level figure. In practice, the framework relies on Librosa and Matplotlib for audio and data visualizations, Modusa for additional MIR-related plots, and ScoreWarp for the warped score that provides the common time axis \cite{mcfeeLibrosaLibrosa01102025,thematplotlibdevelopmentteamMatplotlibVisualizationPython2026,modusa,goeblLetsScoreWarpAgain2025a}.

The key design choice is SVG composition. ScoreWarp already outputs the score as SVG, which preserves the semantic structure of the score while shifting note positions to the performance timeline. LayerIt uses that SVG timeline as the anchor for other visual layers. Raster layers, such as spectrograms, can be embedded when needed, while vector-based layers keep their resolution when the figure is scaled. This makes the framework suitable for publication figures, where reproducibility and visual fidelity both matter.

The system was designed around a small set of practical inputs: a performance audio file, an MEI score, the alignment data that links score elements to time, and optional model output such as beat probabilities or detected beats. Those inputs are transformed into composable elements that can be stacked vertically, reordered, or combined into a single composite visualization. This makes the output more maintainable than a manually edited figure and more expressive than a single-purpose GUI screenshot.

\section{Demo Workflow}

The demo shows how LayerIt supports two related MIR workflows. First, it can generate beat annotation figures in which the score, audio-derived representations, and onset information are aligned in one view. Second, it can generate beat-tracking evaluation figures that expose both the discrete beat detections and the underlying confidence curve, so that model behavior can be inspected beyond a single accuracy score.

In the thesis implementation, the workflow begins with a warped score produced from an MEI file and its alignment map. That score becomes the timeline reference for all additional layers. On top of it, LayerIt can place a spectrogram, beat activation curves, detected beats, downbeats, or other time-based annotations. Because the layer model is modular, the same framework can also accept externally produced SVG or PNG graphics when a workflow already exists in another tool.

This is where the system becomes useful in practice. A beat-tracking result that is otherwise difficult to interpret as a sequence of timestamps can be shown together with the score, the waveform or spectrogram, and the confidence curve. For annotation tasks, the same view makes it easier to see where a human annotation and the model output diverge. For evaluation tasks, it makes the failure mode visible rather than hidden inside a summary metric. That visual context is valuable in the spirit of beat-tracking evaluation work that focuses on correction effort and interpretability, not only on final scores \cite{pintoShiftIfYoua}.

\section{Conclusion}

LayerIt is the result of moving from exploratory thesis prototypes to a focused visualization framework. Its contribution is a reusable way to compose score-aligned MIR figures in Python without depending on manual graphic editing for every new result. The framework is intentionally narrow: it does not replace alignment systems or GUI annotation software, but it provides the visual layer that lets those systems be presented, compared, and extended more effectively.

The current version is already useful for the beat annotation and beat-tracking evaluation scenarios that motivated the thesis, and it also serves as a basis for future MIR visualization work. In particular, the same layer-composition model can be extended to other score-aligned tasks, richer annotation types, or new alignment backends. That makes LayerIt less a single demo than a reusable framework for building clearer MIR figures.

% For bibtex users:
\bibliography{../TeseNando/backmatter/bibliography}

\end{document}
